% 03.sys_env.tex - 第2章 系统开发环境
\section{系统开发环境}

\subsection{硬件环境}
\begin{itemize}[nosep, leftmargin=2em]
    \item CPU: Intel Core i7
    \item 内存: 16GB
    \item 硬盘: 512GB SSD
\end{itemize}

\subsection{软件环境}
\subsubsection{操作系统}
- Windows 11 / Linux (Ubuntu 22.04)

\subsubsection{开发工具}
Visual Studio Code, Git, Docker

\subsubsection{数据库}
MySQL 8.0

\subsection{相关理论基础}
\subsubsection{斯笃克斯定理（Stokes' Theorem）}
斯托克斯定理是微分几何中的重要定理，它将曲面上的曲面积分与该曲面边界的曲线积分联系起来：
\begin{equation}
    \oint_{\partial \Sigma} \mathbf{F} \cdot d\mathbf{r} = \iint_{\Sigma} \nabla \times \mathbf{F} \cdot d\mathbf{S}
    \label{eq:stokes_theorem}
\end{equation}
其中，$\mathbf{F}$ 是向量场，$\partial \Sigma$ 是曲面 $\Sigma$ 的边界曲线。

\subsubsection{质能方程（Mass-Energy Equivalence）}
爱因斯坦提出的质能方程揭示了质量和能量之间的等价关系：
\begin{equation}
    E = mc^2
    \label{eq:mass_energy}
\end{equation}
其中，$E$ 表示能量，$m$ 表示质量，$c$ 表示真空中的光速（$c \approx 3 \times 10^8 \, \text{m/s}$）。

\begin{table}[H]
    \centering
    \caption{系统环境变量配置}
    \label{tbl:env_vars}
    \begin{tabular}{c c c}
        \toprule
        \textbf{环境变量} & \textbf{示例值} & \textbf{描述} \\
        \midrule
        JAVA\_HOME & /usr/lib/jvm/java-11-openjdk & Java 运行环境安装路径 \\
        PATH & \$JAVA\_HOME/bin:\$PATH & 系统可执行文件搜索路径 \\
        MYSQL\_HOST & localhost & MySQL 数据库主机地址 \\
        MYSQL\_PORT & 3306 & MySQL 数据库端口号 \\
        DOCKER\_HOST & unix:///var/run/docker.sock & Docker 守护进程地址 \\
        \bottomrule
    \end{tabular}
\end{table}